\section{The random body}\label{sec:construction}

Fix the numerical constants
\begin{equation}\label{eq:fixed-constants}
 \varepsilon=2^{-30},\qquad a=1+\varepsilon,\qquad
 D=2^{-10}=\varepsilon^{1/3},\qquad v=\frac1{64}.
\end{equation}
For each sufficiently large $n$, put
\begin{equation}\label{eq:radii}
 \eta=n^{-4},\qquad b=a(1+1/n),\qquad
 r_-=a(1-10\log n/n).
\end{equation}
The small window is $B_D$. Its size relative to $a$ will govern the number
of patterns, through
\begin{equation}\label{eq:rates}
 \lambda_0=\log(1+D/a),\qquad
 \gamma=\log\frac{D}{8\sqrt\varepsilon}-2\lambda_0
        =\log4-2\lambda_0.
\end{equation}
Since $\lambda_0<D=2^{-10}$, we have $\gamma>3\lambda_0$.
These constants are fixed; none tends to zero with dimension.

Set
\begin{equation}\label{eq:poisson-measure}
 M=\bigl(\sigma\{u:|u\cdot e_1|>1/a\}\bigr)^{-1},\qquad \mu=M\sigma,
\end{equation}
and let $\Pi$ be a Poisson process on $S^{n-1}$ with intensity $\mu$.
Equivalently, sample a Poisson random integer of mean $M$, then that many
independent $\sigma$-distributed unit normals. Define
\begin{align}
 K&=bB\cap\bigcap_{u\in\Pi}\{x:|u\cdot x|\le1\},\label{eq:body}\\
 K^+&=(b+3\eta)B\cap\bigcap_{u\in\Pi}\{x:|u\cdot x|\le1+3\eta\}.
 \label{eq:buffered-body}
\end{align}
The enlarged body $K^+$ absorbs the displacement of centers during rounding.
Since $M<\infty$, the list of normals is finite almost surely. Both bodies
are compact and centrally symmetric, and $B\subset K$, so $K$ has nonempty
interior.

\begin{lemma}\label{lem:volume}
The body in~\eqref{eq:body} satisfies
\[
 \mathbb P\{|K|>2v\omega_n a^n\}
 \ge \frac{\exp(-1)-2v}{3-2v}>0.
\]
\end{lemma}
\begin{proof}
If $|x|\le a$, the $\mu$-measure of normals with $|u\cdot x|>1$ is at
most one by rotational invariance and~\eqref{eq:poisson-measure}. The
probability that there is no such Poisson normal is therefore at least
$\exp(-1)$. Integrating over $aB$ gives
$\mathbb E|K|\ge\exp(-1)\omega_n a^n$.
On the other hand, $|K|\le\omega_n b^n<3\omega_n a^n$. Thus
$X=|K|/(\omega_n a^n)$ has $0\le X\le3$ and $\mathbb E X\ge\exp(-1)$.
Writing $p=\mathbb P(X>2v)$, we have
$\mathbb E X\le2v(1-p)+3p$, which proves the claim.
\end{proof}

We shall rule out coverings with center intensity at most
\begin{equation}\label{eq:intensity-threshold}
 \rho_0=\frac{R}{2v\omega_n a^n},
\end{equation}
where $R$ is a positive parameter. This threshold does not depend on the
realized volume. Once every such covering is excluded, the definition of
$\theta_{\mathrm{per}}$ gives $\theta_{\mathrm{per}}(K)\ge\rho_0|K|$. The strict volume event in
Lemma~\ref{lem:volume} will then imply $\theta_{\mathrm{per}}(K)>R$, including strictness
at the infimum.
